%% chapters/03_dataset_eda.tex

\section{Dataset and Exploratory Data Analysis}
\label{sec:dataset}

\subsection{IEEE-CIS Fraud Detection Dataset}

The IEEE-CIS Fraud Detection dataset \citep{ieeecisFraud2019} was released
in partnership with Vesta Corporation for a 2019 Kaggle competition.
It comprises two relational tables:

\begin{itemize}
    \item \textbf{Transactions table:} 590,540 rows $\times$ 394 features,
          including transaction amount, product category, card attributes
          (card1--card6), address codes, email domains, distance features
          (dist1, dist2), counting features (C1--C14), timedelta features
          (D1--D15), match features (M1--M9), and 339 Vesta-proprietary
          engineered features (V1--V339).
    \item \textbf{Identity table:} 144,233 rows $\times$ 41 features,
          including device type, device info, and identity features
          (id\_01--id\_38).
\end{itemize}

Tables are joined on \texttt{TransactionID} via a left join, yielding
590,540 merged rows with 431 columns after join.
The binary target label \texttt{isFraud} is provided in the transaction table.

\subsection{Summary Statistics}

\begin{table}[htbp]
  \centering
  \caption{Dataset summary statistics.}
  \label{tab:dataset_stats}
  \begin{tabular}{lr}
    \toprule
    \textbf{Statistic} & \textbf{Value} \\
    \midrule
    Total transactions         & 590,540 \\
    Fraud transactions         & 20,663 \\
    Legitimate transactions    & 569,877 \\
    Fraud prevalence           & 3.499\% \\
    Total features (post-join) & 431 \\
    Numeric features           & 370 \\
    Categorical features       & 61 \\
    Overall null rate          & 47.2\% \\
    Median transaction amount  & \.00 \\
    Median fraud amount        & \.00 \\
    Median legitimate amount   & \.00 \\
    \bottomrule
  \end{tabular}
\end{table}

\subsection{Class Imbalance}

Figure~\ref{fig:class_balance} shows the severe class imbalance:
only 3.5\% of transactions are fraudulent.
This renders raw accuracy meaningless (a classifier predicting all-legitimate
achieves 96.5\% accuracy) and mandates AUPRC as the primary evaluation metric.

\begin{figure}[htbp]
  \centering
  \includegraphics[width=0.45\textwidth]{figures/eda/01_class_balance}
  \caption{Class distribution of the IEEE-CIS dataset.}
  \label{fig:class_balance}
\end{figure}

\subsection{Transaction Amount Distribution}

Fraud transactions exhibit a higher median amount than legitimate ones
(\ vs.\ \), but significant overlap makes amount alone insufficient
for discrimination (Figure~\ref{fig:amount_dist}).
The log-transform normalises the heavy right tail.

\begin{figure}[htbp]
  \centering
  \includegraphics[width=0.65\textwidth]{figures/eda/02_amount_distribution}
  \caption{KDE of $\log(1 + \texttt{TransactionAmt})$ by class.}
  \label{fig:amount_dist}
\end{figure}

\subsection{Temporal Patterns}

The fraud rate varies significantly by hour of day and day of week
(Figure~\ref{fig:temporal_heatmap}), motivating the cyclical time
encoding described in Section~\ref{sec:features}.

\begin{figure}[htbp]
  \centering
  \includegraphics[width=\textwidth]{figures/eda/03_temporal_heatmap}
  \caption{Hourly fraud rate by day of week (0=Monday reference).}
  \label{fig:temporal_heatmap}
\end{figure}

\subsection{Missingness}

The V-feature block (V1--V339) exhibits structured missingness,
with many features exceeding 70\% null rate.
Missingness is non-random (MCAR assumption fails): missing V-features
correlate with device type and, critically, with fraud status.
We exploit this by creating binary missingness indicator features for
all columns exceeding 5\% null rate.
